\chapter{Пространственно-временная классификация временных рядов фМРТ в условиях ограниченной выборки}\label{ch:cls}

В данной главе рассматривается вторая постановка задачи декодирования, дополняющая временное описание пространственным: временные ряды фМРТ одного испытуемого классифицируются по категориям стимулов на основе масок активности классов и проекции ковариационных матриц на касательное пространство риманова многообразия. Размерность признакового описания задается числом активных вокселей и не зависит от объема выборки. Результаты главы опубликованы в работе~\cite{dorin2025enhancing}.

Полный код реализации, использованный для получения результатов главы, доступен в открытом репозитории:
\begin{center}
    {\small \url{https://github.com/DorinDaniil/Spatial-and-Temporal-Characteristics}}
\end{center}

\section{Постановка задачи}\label{sec2:cls}

Временной ряд фМРТ~--- это последовательность из $\mathcal{T}$ сканов:
\begin{equation*}
\mathbf{S} = [\mathbf{F}(1), \ldots, \mathbf{F}(\mathcal{T})] \in \mathcal{X}_s,
\end{equation*}
где $\mathbf{F}(t) \in \mathcal{X}_f$~--- скан в момент $t$. Задана выборка из $N$ рядов одного испытуемого, каждому из которых сопоставлена метка категории стимула:
\begin{equation}\label{eq:cls_dataset}
\mathfrak{D} = \{(\mathbf{S}_i, y_i) \mid i = 1, \ldots, N\}, \qquad y_i \in \{1, \ldots, L\},
\end{equation}
где $L$~--- число категорий стимулов. Требуется построить отображение:
\begin{equation}\label{eq:cls_target}
\mathbf{h}: \mathcal{X}_s \to \{1, \ldots, L\},
\end{equation}
учитывающее пространственно-временные характеристики ряда.

\section{Метод классификации}\label{sec3:cls}

Предлагается методология декодирования с учетом пространственно-временных характеристик (англ. Decoding with Spatio-Temporal characteristics, DeSPoT)~--- комплексный подход для выделения активных областей мозга путем взвешивания вокселей с использованием кросс-корреляции, классификации временных рядов на основе полученных масок активности, а также сегментации временных рядов фМРТ для аугментации. Методология включает алгоритм декодирования активности мозга (англ. Brain Activity Decoder, BAD), модель классификации, использующую этот алгоритм, и метод аугментации. Методология специально разработана для решения задачи декодирования данных одного испытуемого, где объем выборки ограничен. Общая схема методологии приведена на рис.~\ref{fig:cls_overview}.

\begin{figure}[h!]
    \centering
    \includegraphics[width=\textwidth]{cls/Figure_1.pdf}
    \caption{Обзор предложенной методологии. Классифицируются временные ряды фМРТ с учетом индивидуальных структур мозга, не требуя больших объемов данных. Ключевые компоненты включают выделение областей интереса (англ. Region of Interest, ROI) и преобразование пространства признаков с использованием римановой геометрии, что составляет основу извлечения пространственно-временных характеристик. Результатом классификации является распределение вероятностей по категориям стимулов. Кроме того, используются методы аугментации данных для улучшения обучения на данных отдельных испытуемых}
    \label{fig:cls_overview}
\end{figure}

На рис.~\ref{fig:cls_scheme} представлена схема модели классификации, реализующей двухэтапный процесс классификации временных рядов фМРТ одного испытуемого по категориям стимулов. На первом этапе происходит \textit{извлечение масок активности мозга для каждой категории стимулов}, что позволяет выделить области, активируемые в ответ на соответствующие стимулы. На втором этапе выполняется \textit{классификация ряда} с использованием извлеченных масок и отображения на касательное пространство риманова многообразия (англ. Tangent Space Mapping, TSM)~\cite{barachant2011multiclass}.

\begin{figure}[h!]
    \centering
    \includegraphics[width=\textwidth]{figs/cls_scheme.pdf}
    \caption{Схема модели классификации. Модель состоит из двух компонентов: слева~--- извлечение масок активности мозга для каждой категории стимулов, справа~--- классификация на основе полученных масок и TSM. Модель улавливает пространственно-временные характеристики, отражающие степень соответствия объекта и маски активности конкретного класса, и не зависит от размера выборки, что важно при ограниченной выборке одного испытуемого}
    \label{fig:cls_scheme}
\end{figure}

\subsection{Извлечение масок активности мозга}\label{sec3:cls_mask}
Активные области, или области интереса (англ. Region of Interest, ROI), выделяются по корреляции сигнала вокселя со стимулом. Для ускорения каждый скан ряда $\mathbf{S}$ предварительно сжимается 3D average pooling с ядром $k_s$, после чего сигнал вокселя $F_{xyz}(t)$ и бинарный ряд стимуляции $\mathbf{u} = [u(1), \ldots, u(\mathcal{T})]$, $u(t)\in\{0,1\}$, приводятся к нулевому среднему и единичной дисперсии (Z-нормализация). Вычисляется их кросс-корреляция при гемодинамическом лаге $\Delta t$:
\begin{equation}\label{eq:cls_corr}
\rho_{xyz} = \frac{1}{\mathcal{T}-1} \sum_{t} \mathring{u}(t)\, \mathring{F}_{xyz}(t + \Delta t).
\end{equation}
Бинарная маска активности $\mathcal{M} \in \{0,1\}^{X/k_s \times Y/k_s \times Z/k_s}$ задается единицами в $h$ вокселях с наибольшим значением $\rho_{xyz}$.

Формализуем процедуру выделения маски. В основе алгоритма лежат два предположения.

\begin{assumption}\label{as:cls_delay}
Каждый испытуемый имеет индивидуальный постоянный гемодинамический лаг отклика~$\Delta t$, который учитывает время, необходимое для изменения уровня кислорода в крови в ответ на нейронную активность.
\end{assumption}

\begin{assumption}\label{as:cls_peak}
Если область мозга реагирует на стимул, функция кросс-корреляции между вокселями в этой области и временным рядом стимулов достигает пика при гемодинамическом лаге~$\Delta t$. Этот пик указывает на сильную корреляцию между нейронной активностью и стимулом.
\end{assumption}

Предположение~\ref{as:cls_delay} опирается на результат главы~\ref{ch:delay}: гипотеза о наличии гемодинамического лага между предъявлением стимула и BOLD-откликом не была отвергнута~--- ошибка прогноза фМРТ-отклика существенно зависит от лага и на зрительной зоне имеет выраженный минимум (раздел~\ref{sec4:delay_dt}, рис.~\ref{fig:mse_dt}). В настоящую главу переносится сам факт этого влияния, а не конкретное значение лага: $\Delta t$ выступает индивидуальным для испытуемого параметром Алгоритма~\ref{alg:bad}, и кросс-корреляция~\eqref{eq:cls_corr} вычисляется со сдвигом именно на этот лаг.

Алгоритм~\ref{alg:bad} принимает временной ряд фМРТ и бинарный ряд стимулов и имеет три настраиваемых параметра: размер ядра $k_s$, гемодинамический лаг $\Delta t$ и количество активных областей $h$. Частота сканирования $\mu$ используется для перевода лага в число отсчетов, что обеспечивает точный учет временной динамики гемодинамического отклика.

\begin{algorithm}[h!]
    \caption{Декодирование активности мозга (англ. Brain Activity Decoder, BAD)}\label{alg:bad}
    \KwIn{временной ряд фМРТ $\mathbf{S}$; бинарный ряд стимулов $\mathbf{u}$; параметры $k_s$, $\Delta t$, $h$}
    \KwOut{маска активности $\mathcal{M}$ и ее версия $\mathcal{M}^{\uparrow}$ в исходном разрешении}
    \nl \textbf{Сжатие}: 3D average pooling с ядром $k_s$:
    $\mathbf{S} \rightarrow \mathbf{V}$, $\mathbf{V} \in \mathbb{R}^{\mathcal{T} \times X/k_s \times Y/k_s \times Z/k_s}$\;
    \nl \textbf{Z-нормализация}: нормализация временных рядов вокселей и ряда стимулов:
    $\mathbf{V} \rightarrow \mathring{\mathbf{V}}$, $\mathbf{u} \rightarrow \mathring{\mathbf{u}}$\;
    \nl \textbf{Кросс-корреляция}: для каждого вокселя $(x, y, z)$ вычисляется
    \[
    \rho_{xyz}(p) = \frac{1}{\mathcal{T}-1} \sum_{t=1}^{\mathcal{T}-p} \mathring{u}(t)\, \mathring{V}_{xyz}(t+p), \quad p = 0, \ldots, \mathcal{T}-1\;
    \]
    \nl \textbf{Определение маски}: вычисляется $p_{\Delta t} = \lfloor \mu \Delta t \rfloor$; создается бинарная маска $\mathcal{M} \in \{0,1\}^{X/k_s \times Y/k_s \times Z/k_s}$ с единицами в позициях $h$ наибольших значений $\rho_{xyz}(p_{\Delta t})$\;
    \nl \textbf{Увеличение разрешения}: для получения $\mathcal{M}^{\uparrow} \in \{0,1\}^{X \times Y \times Z}$ каждому элементу присваивается значение ближайшего элемента $\mathcal{M}$:
    \[
    \mathcal{M}^{\uparrow}(x, y, z) = \mathcal{M}\left(\left\lfloor \tfrac{x}{k_s} \right\rfloor, \left\lfloor \tfrac{y}{k_s} \right\rfloor, \left\lfloor \tfrac{z}{k_s} \right\rfloor\right)\;
    \]
    \Return{$\mathcal{M}$, $\mathcal{M}^{\uparrow}$}
\end{algorithm}

Кросс-корреляция~\eqref{eq:cls_corr} есть значение функции $\rho_{xyz}(p)$ из Алгоритма~\ref{alg:bad} в точке $p = p_{\Delta t}$. Для проверки статистической значимости областей, полученных Алгоритмом~\ref{alg:bad}, проверяется значимость корреляции между временным рядом стимулов и временным рядом каждого вокселя. Обозначим $\mathbf{v}$~--- временной ряд вокселя и $\tau_{\mathbf{v}\mathbf{u}}$~--- коэффициент ранговой корреляции Кендалла. Применяется следующий критерий:
\begin{center}
    \begin{tabular}{rl}
    \textit{Выборки:}                        & $\mathbf{v}=\left[v(1),\ldots,v(\mathcal{T})\right]$\\
                                    & $\mathbf{u}=\left[u(1),\ldots,u(\mathcal{T})\right]$\\
                                    & выборки связаны\\
    \textit{Нулевая гипотеза:}                & $H_0\colon \tau_{\mathbf{v}\mathbf{u}}=0$ \\
    \textit{Альтернативная гипотеза:}         & $H_1\colon \tau_{\mathbf{v}\mathbf{u}}>0$ \\
    \textit{Статистика критерия:}             & $\hat{\tau}_{\mathbf{v}\mathbf{u}}$ \\
    \textit{Распределение при нулевой гипотезе:} & табличное\\
    \end{tabular}
\end{center}
Критерий применяется к каждому вокселю, что образует задачу множественной проверки гипотез. Для контроля групповой вероятности ошибки первого рода (англ. Family-Wise Error Rate, FWER) на заданном уровне значимости используется поправка Холма. Таким образом отсеиваются ложные срабатывания и выделяются воксели со статистически значимой корреляцией.

Маска строится отдельно для каждой категории $y = 1, \ldots, L$. Применение Алгоритма~\ref{alg:bad} к отдельным объектам класса не гарантирует, что полученная маска соответствует области мозга, ответственной за стимуляцию класса $y$, из-за ограничений процедуры фМРТ, описанных в главе~\ref{sec1}. Поэтому для повышения статистической значимости ряды и стимулы объектов класса $y$ объединяются по времени:
\begin{align*}
& \widetilde{\mathbf{S}}^{y} = \mathbf{S}_1 \oplus \mathbf{S}_2 \oplus \ldots \oplus \mathbf{S}_{N_y}, \quad \widetilde{\mathbf{S}}^{y} \in \mathbb{R}^{\mathcal{T} N_y \times X \times Y \times Z},\\
& \widetilde{\mathbf{u}}^{y} = \mathbf{u}_1 \oplus \mathbf{u}_2 \oplus \ldots \oplus \mathbf{u}_{N_y}, \quad \widetilde{\mathbf{u}}^{y} \in \{0, 1\}^{\mathcal{T} N_y},
\end{align*}
где $N_y$~--- число объектов класса $y$, а $\oplus$ обозначает конкатенацию по временной оси. К конкатенированным данным применяется Алгоритм~\ref{alg:bad} с параметрами $h_y$, $k_s$, $\Delta t$; шаг увеличения разрешения опускается, поскольку применение average pooling уменьшает влияние аномальных и зашумленных вокселей. Результат~--- маска $\mathcal{M}^{y} \in \{0,1\}^{X/k_s \times Y/k_s \times Z/k_s}$. Усреднение по объектам класса повышает статистическую значимость маски и ее устойчивость к шуму отдельных записей.

Формально определим результат применения маски к временному ряду.

\begin{definition}\label{def:cls_masking}
Результат применения маски активности $\mathcal{M}$ к ряду $\mathbf{S} = [\mathbf{F}(1), \ldots, \mathbf{F}(\mathcal{T})]$ определяется как матрица
\begin{equation*}
\mathbf{Y} = \mathcal{M} \odot \mathbf{S} \in \mathbb{R}^{h \times \mathcal{T}},
\end{equation*}
где $\odot$ обозначает произведение Адамара, а столбец $t$ матрицы $\mathbf{Y}$ состоит из значений скана $\mathbf{F}(t)$, соответствующих ненулевым элементам маски, $h$~--- число таких элементов.
\end{definition}

\subsection{Классификация временных рядов}\label{sec3:cls_clf}
Задан набор данных $\mathfrak{D}$, см.~\eqref{eq:cls_dataset}, и маски активности $\mathcal{M}^{y}$, $y = 1, \ldots, L$. Искомое отображение $\mathbf{h}: \mathcal{X}_s \to \{1,\ldots,L\}$ представляется суперпозицией
\begin{equation}\label{eq:cls_model}
\mathbf{h} = \boldsymbol{\eta} \circ \boldsymbol{\psi} \circ \mathcal{A},
\end{equation}
где $\mathcal{A}$~--- 3D average pooling с ядром $k_s$, $\boldsymbol{\psi}$~--- кодировщик пространственно-временных характеристик, $\boldsymbol{\eta}: \mathbb{R}^{d} \to \{1,\ldots,L\}$~--- классификатор.

Кодировщик является конкатенацией отображений для каждой категории стимула:
\begin{equation}\label{eq:cls_psi}
\boldsymbol{\psi} = \boldsymbol{\psi}_1 \oplus \cdots \oplus \boldsymbol{\psi}_L, \qquad \boldsymbol{\psi}_y = \boldsymbol{\pi}_y \circ \mathbf{m}_y, \qquad d = \sum_{y=1}^{L} d_y,
\end{equation}
где $\mathbf{m}_y$ применяет маску $\mathcal{M}^{y}$ в смысле Определения~\ref{def:cls_masking} и отбирает $h_y$ активных вокселей, формируя матрицу из $\mathbb{R}^{h_y \times \mathcal{T}}$, а $\boldsymbol{\pi}_y$ проецирует ее на касательное пространство риманова многообразия с последующей векторизацией. Размерность определяется в разделе~\ref{sec3:cls_tsm}:
\begin{equation}\label{eq:cls_dim}
d_y = \frac{h_y(h_y+1)}{2}.
\end{equation}
Ключевым компонентом модели является кодировщик $\boldsymbol{\psi}$, который транслирует временные ряды в евклидово пространство, проецируя их на касательное пространство риманова многообразия, которое само по себе является евклидовым. Это преобразование упрощает последующую задачу классификации, эффективно извлекая пространственно-временные характеристики, демонстрирующие степень соответствия между объектом и маской активности конкретного класса. В результате выбор классификатора становится простым, и в данной работе в качестве $\boldsymbol{\eta}$ рассматриваются логистическая регрессия и двухслойный перцептрон.

\subsection{Проекция на касательное пространство риманова многообразия}\label{sec3:cls_tsm}
Для упрощения изложения и без потери общности индекс класса $y$ опускается во всех обозначениях. Отображение $\boldsymbol{\pi}: \mathbb{R}^{h \times \mathcal{T}} \to \mathbb{R}^{d}$ называется отображением на касательное пространство (англ. Tangent Space Mapping, TSM)~\cite{barachant2011multiclass}. Оно преобразует пространство объектов в терминах римановой геометрии~\cite{barachant2010riemannian, barachant2011multiclass}: выборка ковариационных матриц объектов проецируется с риманова многообразия, в котором скалярное произведение не определено, на касательную плоскость, которая является евклидовым пространством. Первым шагом метода является формирование пространства центрированных признаков. Для $i$-го объекта матрица активных вокселей $\mathbf{Y}_i = \mathcal{M} \odot \mathbf{S}_i$ из Определения~\ref{def:cls_masking} центрируется по времени:
\begin{equation*}
\mathbf{Y}_i = \left[
\begin{array}{cccc}
y^{i}_{11} & y^{i}_{12} & \ldots & y^{i}_{1\mathcal{T}}\\
y^{i}_{21} & y^{i}_{22} & \ldots & y^{i}_{2\mathcal{T}}\\
\vdots & \vdots & \ddots & \vdots\\
y^{i}_{h1} & y^{i}_{h2} & \ldots & y^{i}_{h\mathcal{T}}
\end{array}
\right]
= \begin{bmatrix}
\mathbf{y}^{i}_1\\
\mathbf{y}^{i}_2\\
\vdots\\
\mathbf{y}^{i}_{h}
\end{bmatrix} \in \mathbb{R}^{h \times \mathcal{T}}, \qquad i = 1, \ldots, N,
\end{equation*}
где $\mathbf{y}^{i}_k$, $k = 1, \ldots, h$,~--- центрированные временные ряды активных вокселей. Затем вычисляется ковариационная матрица многомерного временного ряда $\mathbf{Y}_i$:
\begin{equation}\label{eq:cls_cov}
\mathbf{R}_i = \frac{1}{\mathcal{T}-1}\, \mathbf{Y}_i \mathbf{Y}_i^{\top} \in \mathbb{S}_{+}^{h}, \qquad i = 1, \ldots, N,
\end{equation}
где $\mathbb{S}_{+}^{h}$~--- пространство симметричных положительно определенных матриц размера $h \times h$. Известно, что это пространство образует риманово многообразие~$\mathcal{R}$~\cite{barachant2010riemannian}.

В каждой точке риманова многообразия существует касательная плоскость с определенным на ней скалярным произведением. Общая касательная плоскость для проецирования всех ковариационных матриц выборки формируется в точке геометрического среднего по римановой метрике известных ковариационных матриц. Геометрическое среднее SPD-матриц задается как
\begin{equation}\label{eq:cls_mean}
\bar{\mathbf{R}} = \mathfrak{G}\left(\mathbf{R}_1, \ldots, \mathbf{R}_N\right) = \operatorname*{arg\,min}_{\mathbf{R} \in \mathbb{S}_{+}^{h}} \sum_{i=1}^{N} \delta_{\mathcal{R}}^2(\mathbf{R}, \mathbf{R}_i),
\end{equation}
где риманова метрика, или геодезическое расстояние, от $\mathbf{R}_1$ до $\mathbf{R}_2$ равна
\begin{equation}\label{eq:cls_dist}
\delta_{\mathcal{R}}(\mathbf{R}_1, \mathbf{R}_2) = \big\| \log\big(\mathbf{R}_1^{-1}\mathbf{R}_2\big) \big\|_F = \Big( \sum_{l=1}^{h} \log^2 \lambda_l \Big)^{1/2},
\end{equation}
а $\lambda_l$~--- собственные значения матрицы $\mathbf{R}_1^{-1}\mathbf{R}_2$. Существование и единственность среднего~\eqref{eq:cls_mean} обсуждаются в разделе~\ref{sec3:cls_theory}.

Для каждой ковариационной матрицы $\mathbf{R}_i$ существует ее проекция $\mathbf{P}_i$ на касательное пространство $T_{\bar{\mathbf{R}}}$ в точке $\bar{\mathbf{R}}$. Проекция задается логарифмическим отображением
\begin{equation}\label{eq:cls_log}
\mathbf{P}_i = \operatorname{Log}_{\bar{\mathbf{R}}}(\mathbf{R}_i) = \bar{\mathbf{R}}^{1/2} \log\!\big( \bar{\mathbf{R}}^{-1/2} \mathbf{R}_i \bar{\mathbf{R}}^{-1/2} \big) \bar{\mathbf{R}}^{1/2}, \qquad \mathbf{P}_i \in \mathbb{R}^{h \times h},
\end{equation}
а обратное отображение с касательного пространства на многообразие~--- экспоненциальным:
\begin{equation}\label{eq:cls_exp}
\operatorname{Exp}_{\bar{\mathbf{R}}}(\mathbf{P}_i) = \mathbf{R}_i = \bar{\mathbf{R}}^{1/2} \exp\!\big( \bar{\mathbf{R}}^{-1/2} \mathbf{P}_i \bar{\mathbf{R}}^{-1/2} \big) \bar{\mathbf{R}}^{1/2}.
\end{equation}
Эти операции вычисляются через спектральное разложение матриц. Для матрицы $\mathbf{Q} = \mathbf{U} \mathbf{D} \mathbf{U}^{-1}$, где $\mathbf{D}$~--- диагональная матрица собственных значений, а $\mathbf{U}$~--- матрица собственных векторов,
\begin{equation}\label{eq:cls_matfun}
\mathbf{Q}^{\pm 1/2} = \mathbf{U} \mathbf{D}^{\pm 1/2} \mathbf{U}^{-1}, \qquad \log(\mathbf{Q}) = \mathbf{U} \log(\mathbf{D}) \mathbf{U}^{-1}, \qquad \exp(\mathbf{Q}) = \mathbf{U} \exp(\mathbf{D}) \mathbf{U}^{-1}.
\end{equation}

После проекции каждой ковариационной матрицы $\mathbf{R}_i$ с использованием $\operatorname{Log}_{\bar{\mathbf{R}}}(\,\cdot\,)$ получаются симметричные матрицы $\mathbf{P}_i \in \mathbb{R}^{h \times h}$. Последним шагом метода является векторизация. Длина касательного вектора в точке $\bar{\mathbf{R}}$ измеряется не нормой Фробениуса самой матрицы $\mathbf{P}_i$, а нормой ее отбеленной формы
\begin{equation}\label{eq:cls_white}
\mathbf{S}_i = \bar{\mathbf{R}}^{-1/2} \mathbf{P}_i \bar{\mathbf{R}}^{-1/2} = \log\!\big( \bar{\mathbf{R}}^{-1/2} \mathbf{R}_i \bar{\mathbf{R}}^{-1/2} \big),
\end{equation}
см. раздел~\ref{sec3:cls_theory}. Поэтому векторизуется матрица $\mathbf{S}_i$: элементы ее верхнего треугольника последовательно записываются в вектор, причем диагональные элементы берутся с коэффициентом~$1$, а внедиагональные~--- с коэффициентом~$\sqrt{2}$. В результате получаются эмбеддинги
\begin{equation}\label{eq:cls_emb}
\mathbf{p}_i = \operatorname{vec}_{\sqrt{2}}(\mathbf{S}_i) = \begin{bmatrix}
S_{11} & \sqrt{2}\, S_{12} & \ldots & \sqrt{2}\, S_{1h} & S_{22} & \sqrt{2}\, S_{23} & \ldots & S_{hh}
\end{bmatrix}^{\top} \in \mathbb{R}^{d},
\end{equation}
размерности $d = h(h+1)/2$. Коэффициент $\sqrt{2}$ согласует евклидову норму вектора с нормой Фробениуса матрицы: каждый внедиагональный элемент симметричной матрицы входит в $\|\mathbf{S}_i\|_F^2 = \sum_k S_{kk}^2 + 2\sum_{k<l} S_{kl}^2$ дважды, а в верхний треугольник~--- один раз, и множитель $\sqrt{2}$ восполняет недостающую двойку, так что $\|\mathbf{p}_i\|_2 = \|\mathbf{S}_i\|_F$. Без этого коэффициента евклидовы расстояния между эмбеддингами занижали бы вклад внедиагональных элементов, а $L_2$-регуляризация классификатора штрафовала бы диагональные и внедиагональные направления по-разному. Тем самым $\mathbf{p}_i$~--- координаты касательного вектора $\mathbf{P}_i$ в ортонормированном базисе касательного пространства; их метрические свойства доказываются в разделе~\ref{sec3:cls_theory}. Процесс проекции ковариационных матриц на касательное пространство, а затем обратно на многообразие проиллюстрирован на рис.~\ref{fig:rieman}.

\begin{figure}[h!]
    \centering
    \includegraphics[width=0.7\textwidth]{cls/Figure_3.pdf}
    \caption{Схема операций римановой геометрии. Ковариационная матрица $\mathbf{R}_j$ проецируется на касательное пространство $T_\mathbf{R}$ в точке $\mathbf{R}$ (в предложенном методе $\mathbf{R} = \bar{\mathbf{R}}$) с использованием логарифмического отображения $\operatorname{Log}_{\mathbf{R}}(\mathbf{R}_j)$. Полученная проекция $\mathbf{P}_j$ затем отображается обратно на риманово многообразие $\mathcal{R}$ с использованием экспоненциального отображения $\operatorname{Exp}_{\mathbf{R}}(\mathbf{P}_j)$. $\delta_{\mathcal{R}}(\mathbf{R},\ \mathbf{R}_j)$~--- геодезическое расстояние между $\mathbf{R}$ и $\mathbf{R}_j$.}
    \label{fig:rieman}
\end{figure}

\subsection{Сегментация временных рядов для аугментации}\label{sec3:cls_aug}

Записи фМРТ обычно содержат одно протяженное измерение для каждого испытуемого, тогда как для обучения модели~\eqref{eq:cls_model} требуется выборка размеченных рядов фиксированной длины. Для формирования такой выборки предлагается метод сегментации временного ряда, выделяющий сегменты с преобладающей категорией стимулов.

\begin{definition}
Фрагмент временного ряда $\mathbf{S}_{\mathrm{rec}} = [\mathbf{F}(1), \ldots, \mathbf{F}(\mathcal{T}_{\mathrm{rec}})]$~--- любая его подпоследовательность
\begin{equation*}
   [\mathbf{F}(t_1), \ldots, \mathbf{F}(t_k)], \quad 1 \leqslant t_1 < \ldots < t_k \leqslant \mathcal{T}_{\mathrm{rec}}.
\end{equation*}
\end{definition}

\begin{definition}
Сегмент временного ряда $\mathbf{S}_{\mathrm{rec}}$~--- его непрерывный фрагмент
\begin{equation*}
[\mathbf{F}(t)]_{t=t_0}^{t_1}, \quad 1 \leqslant t_0 < t_1 \leqslant \mathcal{T}_{\mathrm{rec}}.
\end{equation*}
\end{definition}

\begin{definition}
Сегментация временного ряда~--- отображение $\mathcal{G}$, которое ставит в соответствие ряду $\mathbf{S}_{\mathrm{rec}}$ множество его сегментов.
\end{definition}

Пусть задана полная запись $\mathbf{S}_{\mathrm{rec}} = [\mathbf{F}(1), \ldots, \mathbf{F}(\mathcal{T}_{\mathrm{rec}})]$ длиной $\mathcal{T}_{\mathrm{rec}}$ и категориальный ряд стимулов
\begin{equation}\label{eq:cls_stimseq}
	\mathbf{u}_{\mathrm{rec}} = [u_{\mathrm{rec}}(1), \ldots, u_{\mathrm{rec}}(\mathcal{T}_{\mathrm{rec}})], \quad u_{\mathrm{rec}}(t) \in \{0, 1, \dots, L\},
\end{equation}
где класс $0$ обозначает периоды отдыха испытуемого. Цель~--- сегментировать запись на интервалы фиксированной длины $\mathcal{T}$ так, чтобы каждый сегмент преимущественно содержал одну категорию стимулов.

\begin{assumption}\label{as:cls_rest}
Испытуемый имел достаточно времени для отдыха между предъявлением стимулов из разных категорий: существует $\mathcal{T} > 0$ такое, что для любого сегмента длины $\mathcal{T}$ суммарная длина моментов стимуляции в нем меньше $\mathcal{T}$.
\end{assumption}

Фиксируем $\mathcal{T}$ в соответствии с Предположением~\ref{as:cls_rest}. Для каждой категории $y \in \{1,\dots,L\}$ на первом шаге находятся начальные и конечные индексы интервалов стимуляции класса $y$, упорядоченные по возрастанию:
\begin{align*}
\{b_{1}, \ldots, b_{N_{y}}\} &= \{t \mid u_{\mathrm{rec}}(t)=y,\ u_{\mathrm{rec}}(t-1)=0\},\\
\{e_{1}, \ldots, e_{N_{y}}\} &= \{t \mid u_{\mathrm{rec}}(t)=y,\ u_{\mathrm{rec}}(t+1)=0\},
\end{align*}
где $N_{y}$~--- число интервалов категории $y$. На втором шаге каждый интервал $[b_{j}, e_{j}]$ длины $\mathcal{T}_{j} = e_{j} - b_{j} + 1$ дополняется периодами отдыха слева и справа до сегмента длины $\mathcal{T}$:
\begin{equation*}
\mathbf{S}_{j} = [\mathbf{F}(b_{j}-r_1), \ldots, \mathbf{F}(e_{j}+r_2)],
\qquad
r_1 = \lfloor (\mathcal{T} - \mathcal{T}_{j})/2 \rfloor, \quad r_2 = \lceil (\mathcal{T} - \mathcal{T}_{j})/2 \rceil,
\end{equation*}
где $j = 1,\ldots, N_{y}$, а $\lfloor \cdot \rfloor$ и $\lceil \cdot \rceil$ обозначают округление вниз и вверх. Аналогично сегментируется ряд стимулов; бинарный ряд стимуляции $\mathbf{u}$ объекта класса $y$ задается индикатором $u(t) = \mathds{1}\{u_{\mathrm{rec}}(t) = y\}$. Объединяя все категории, получаем выборку из $N = \sum_{y=1}^{L} N_y$ сегментов длины $\mathcal{T}$ с метками классов.

\section{Теоретические свойства проекции на касательное пространство}\label{sec3:cls_theory}

Кодировщик~$\boldsymbol{\psi}$ из~\eqref{eq:cls_psi} заменяет ковариационную матрицу объекта ее проекцией на касательное пространство в точке риманова среднего, после чего работает линейный классификатор. В настоящем разделе доказываются три свойства этой конструкции. Во-первых, признаки не зависят от единиц измерения сигнала и от невырожденного смешивания отобранных вокселей. Во-вторых, евклидовы расстояния между признаками отличаются от геодезических расстояний между ковариациями не более чем в контролируемое число раз, зависящее только от радиуса облака объектов. В-третьих, ошибка оценивания ковариации по $\mathcal{T}$ сканам не зависит от истинной ковариации и убывает как $h/\sqrt{\mathcal{T}}$, что задает ограничение на число отбираемых вокселей~$h$.

\textbf{Риманова структура на множестве ковариаций.} Множество $\mathbb{S}_{+}^{h}$ снабжается аффинно-инвариантной римановой метрикой~\cite{bhatia2007positive}: касательное пространство $T_{\mathbf{R}}$ в точке $\mathbf{R} \in \mathbb{S}_{+}^{h}$ есть пространство симметричных матриц $\mathrm{Sym}(h)$, а скалярное произведение касательных векторов $\mathbf{X}, \mathbf{Y} \in \mathrm{Sym}(h)$ задается формулой
\begin{equation}\label{eq:cls_metric}
\langle \mathbf{X}, \mathbf{Y} \rangle_{\mathbf{R}} = \operatorname{tr}\big(\mathbf{R}^{-1} \mathbf{X} \mathbf{R}^{-1} \mathbf{Y}\big),
\qquad
\|\mathbf{X}\|_{\mathbf{R}} = \big\| \mathbf{R}^{-1/2} \mathbf{X} \mathbf{R}^{-1/2} \big\|_F .
\end{equation}
Геодезическое расстояние~\eqref{eq:cls_dist} допускает эквивалентную симметричную запись
\begin{equation}\label{eq:cls_dist_log}
\delta_{\mathcal{R}}(\mathbf{R}_1, \mathbf{R}_2) = \big\| \log\big( \mathbf{R}_1^{-1/2} \mathbf{R}_2 \mathbf{R}_1^{-1/2} \big) \big\|_F ,
\end{equation}
поскольку матрицы $\mathbf{R}_1^{-1}\mathbf{R}_2$ и $\mathbf{R}_1^{-1/2}\mathbf{R}_2\mathbf{R}_1^{-1/2}$ подобны и имеют одни и те же собственные значения $\lambda_l > 0$, а норма Фробениуса симметричной матрицы равна корню из суммы квадратов ее собственных значений; в отличие от $\log(\mathbf{R}_1^{-1}\mathbf{R}_2)$, матрица под нормой в~\eqref{eq:cls_dist_log} симметрична.

\begin{lemma}\label{lem:cls_inv}
Для любой обратимой матрицы $\mathbf{A} \in \mathbb{R}^{h \times h}$ и любых $\mathbf{R}_1, \mathbf{R}_2 \in \mathbb{S}_{+}^{h}$
\begin{equation}\label{eq:cls_inv}
\delta_{\mathcal{R}}\big(\mathbf{A}\mathbf{R}_1\mathbf{A}^{\top}, \mathbf{A}\mathbf{R}_2\mathbf{A}^{\top}\big) = \delta_{\mathcal{R}}(\mathbf{R}_1, \mathbf{R}_2).
\end{equation}
\end{lemma}
\begin{proof}
Расстояние зависит только от собственных значений матрицы $\mathbf{R}_1^{-1}\mathbf{R}_2$. Матрица $(\mathbf{A}\mathbf{R}_1\mathbf{A}^{\top})^{-1}\mathbf{A}\mathbf{R}_2\mathbf{A}^{\top} = \mathbf{A}^{-\top}\mathbf{R}_1^{-1}\mathbf{R}_2\mathbf{A}^{\top}$ подобна $\mathbf{R}_1^{-1}\mathbf{R}_2$, а подобные матрицы имеют одинаковые собственные значения.
\end{proof}

Риманово среднее $\bar{\mathbf{R}}$ из раздела~\ref{sec3:cls_tsm} существует и единственно, поскольку функция $\mathbf{R} \mapsto \delta_{\mathcal{R}}^2(\mathbf{R}, \mathbf{R}_i)$ строго геодезически выпукла на многообразии неположительной кривизны~\cite{bhatia2007positive, moakher2005differential}. Оно характеризуется уравнением
\begin{equation}\label{eq:cls_karcher}
\sum_{i=1}^{N} \log\big( \bar{\mathbf{R}}^{-1/2} \mathbf{R}_i \bar{\mathbf{R}}^{-1/2} \big) = \mathbf{0}.
\end{equation}

\textbf{Координаты в касательном пространстве.} Проекция $\mathbf{P}_i = \operatorname{Log}_{\bar{\mathbf{R}}}(\mathbf{R}_i)$ является касательным вектором в точке $\bar{\mathbf{R}}$, и ее длина измеряется нормой~\eqref{eq:cls_metric}. Напомним обозначения раздела~\ref{sec3:cls_tsm}, см.~\eqref{eq:cls_white} и~\eqref{eq:cls_emb}: отбеленная форма касательного вектора и признак объекта равны
\begin{equation}\label{eq:cls_feat}
\mathbf{S}_i = \bar{\mathbf{R}}^{-1/2} \mathbf{P}_i \bar{\mathbf{R}}^{-1/2} = \log\big( \bar{\mathbf{R}}^{-1/2} \mathbf{R}_i \bar{\mathbf{R}}^{-1/2} \big),
\qquad
\mathbf{p}_i = \operatorname{vec}_{\sqrt{2}}(\mathbf{S}_i) \in \mathbb{R}^{d},
\end{equation}
где $\operatorname{vec}_{\sqrt{2}}$~--- векторизация~\eqref{eq:cls_emb} и $d = h(h+1)/2$.

\begin{lemma}\label{lem:cls_vec}
Отображение $\operatorname{vec}_{\sqrt{2}}$ является линейной изометрией пространств $(\mathrm{Sym}(h), \|\cdot\|_F)$ и $(\mathbb{R}^{d}, \|\cdot\|_2)$. Как следствие,
\begin{equation}\label{eq:cls_vec_iso}
\|\mathbf{p}_i\|_2 = \|\mathbf{P}_i\|_{\bar{\mathbf{R}}},
\qquad
\|\mathbf{p}_i - \mathbf{p}_j\|_2 = \|\mathbf{P}_i - \mathbf{P}_j\|_{\bar{\mathbf{R}}} .
\end{equation}
\end{lemma}
\begin{proof}
Для симметричной $\mathbf{S}$ каждый внедиагональный элемент входит в норму Фробениуса дважды:
$\|\mathbf{S}\|_F^2 = \sum_{k} S_{kk}^2 + 2\sum_{k<l} S_{kl}^2 = \|\operatorname{vec}_{\sqrt{2}}(\mathbf{S})\|_2^2$. Линейность очевидна. Второе утверждение следует из~\eqref{eq:cls_metric} и~\eqref{eq:cls_feat}.
\end{proof}

\begin{remark}
Из уравнения~\eqref{eq:cls_karcher} и линейности $\operatorname{vec}_{\sqrt{2}}$ следует $\sum_i \mathbf{p}_i = \mathbf{0}$: признаки обучающей выборки центрированы автоматически.
\end{remark}

Далее используются два элементарных факта о матричных функциях. Для симметричной $\mathbf{M} = \mathbf{U} \operatorname{diag}(\lambda_1, \ldots, \lambda_h) \mathbf{U}^{\top}$ с ортогональной $\mathbf{U}$ и функции $f$ положено $f(\mathbf{M}) = \mathbf{U} \operatorname{diag}(f(\lambda_1), \ldots, f(\lambda_h)) \mathbf{U}^{\top}$, что согласуется с~\eqref{eq:cls_matfun}. Отображение $\exp\colon \mathrm{Sym}(h) \to \mathbb{S}_{+}^{h}$ есть диффеоморфизм с обратным~$\log$.

\begin{lemma}\label{lem:cls_conj}
Для ортогональной $\mathbf{Q}$ и симметричной $\mathbf{M}$ выполнено $f(\mathbf{Q}\mathbf{M}\mathbf{Q}^{\top}) = \mathbf{Q} f(\mathbf{M}) \mathbf{Q}^{\top}$.
\end{lemma}
\begin{proof}
Если $\mathbf{M} = \mathbf{U}\operatorname{diag}(\lambda)\mathbf{U}^{\top}$, то $\mathbf{Q}\mathbf{M}\mathbf{Q}^{\top} = (\mathbf{Q}\mathbf{U})\operatorname{diag}(\lambda)(\mathbf{Q}\mathbf{U})^{\top}$~--- спектральное разложение с теми же собственными значениями и ортогональной матрицей $\mathbf{Q}\mathbf{U}$.
\end{proof}

\begin{lemma}\label{lem:cls_duhamel}
Для симметричных $\mathbf{S}$ и $\mathbf{X}$ дифференциал матричной экспоненты равен
\begin{equation}\label{eq:cls_duhamel}
d\exp_{\mathbf{S}}[\mathbf{X}] = \frac{d}{dt} e^{\mathbf{S} + t\mathbf{X}} \Big|_{t=0} = \int_{0}^{1} e^{(1-u)\mathbf{S}} \mathbf{X} e^{u\mathbf{S}} \, du .
\end{equation}
\end{lemma}
\begin{proof}
Положим $\mathbf{G}(u) = e^{(1-u)\mathbf{S}} e^{u(\mathbf{S}+t\mathbf{X})}$. Тогда $\mathbf{G}(1) - \mathbf{G}(0) = e^{\mathbf{S}+t\mathbf{X}} - e^{\mathbf{S}}$, а производная
$\mathbf{G}'(u) = e^{(1-u)\mathbf{S}} \big( -\mathbf{S} + \mathbf{S} + t\mathbf{X} \big) e^{u(\mathbf{S}+t\mathbf{X})} = t\, e^{(1-u)\mathbf{S}} \mathbf{X} e^{u(\mathbf{S}+t\mathbf{X})}$.
Интегрируя по $u \in [0,1]$, деля на $t$ и переходя к пределу $t \to 0$, получаем~\eqref{eq:cls_duhamel}.
\end{proof}

\textbf{Инвариантность к преобразованию каналов.} Первый вопрос к методу, работающему с многоканальной записью: зависит ли ответ от единиц измерения каналов и от их смешивания. Следующее предложение показывает, что для описанной конструкции ответ отрицательный.

\begin{proposition}\label{prop:cls_affine}
Пусть $\mathbf{A} \in \mathbb{R}^{h \times h}$ обратима и все ковариационные матрицы преобразованы как $\mathbf{R}_i' = \mathbf{A} \mathbf{R}_i \mathbf{A}^{\top}$, $i = 1, \ldots, N$, что отвечает замене рядов $\mathbf{Y}_i$ на $\mathbf{A}\mathbf{Y}_i$. Тогда
\begin{enumerate}
    \item[(а)] риманово среднее преобразуется согласованно: $\bar{\mathbf{R}}' = \mathbf{A} \bar{\mathbf{R}} \mathbf{A}^{\top}$;
    \item[(б)] существует ортогональная матрица $\mathbf{O} \in \mathbb{R}^{d \times d}$, одна и та же для всех объектов, такая что $\mathbf{p}_i' = \mathbf{O} \mathbf{p}_i$;
    \item[(в)] решение $L_2$-регуляризованной логистической регрессии на признаках $\{\mathbf{p}_i'\}$ равно $\hat{\mathbf{W}}' = \hat{\mathbf{W}} \mathbf{O}^{\top}$, где $\hat{\mathbf{W}}$~--- решение на $\{\mathbf{p}_i\}$, и предсказания на любом объекте совпадают.
\end{enumerate}
\end{proposition}
\begin{proof}
(а) По лемме~\ref{lem:cls_inv} для любой $\mathbf{R} \in \mathbb{S}_{+}^{h}$ выполнено $\sum_i \delta_{\mathcal{R}}^2(\mathbf{A}\mathbf{R}\mathbf{A}^{\top}, \mathbf{R}_i') = \sum_i \delta_{\mathcal{R}}^2(\mathbf{R}, \mathbf{R}_i)$. Отображение $\mathbf{R} \mapsto \mathbf{A}\mathbf{R}\mathbf{A}^{\top}$~--- биекция $\mathbb{S}_{+}^{h}$ на себя, переводящая целевую функцию задачи о среднем для исходных матриц в целевую функцию задачи для преобразованных. Поэтому оно переводит минимизатор в минимизатор, и по единственности $\bar{\mathbf{R}}' = \mathbf{A}\bar{\mathbf{R}}\mathbf{A}^{\top}$.

(б) Положим $\mathbf{Q} = (\mathbf{A}\bar{\mathbf{R}}\mathbf{A}^{\top})^{-1/2} \mathbf{A} \bar{\mathbf{R}}^{1/2}$. Прямая проверка дает
\begin{equation*}
\mathbf{Q}\mathbf{Q}^{\top} = (\mathbf{A}\bar{\mathbf{R}}\mathbf{A}^{\top})^{-1/2} \mathbf{A}\bar{\mathbf{R}}\mathbf{A}^{\top} (\mathbf{A}\bar{\mathbf{R}}\mathbf{A}^{\top})^{-1/2} = \mathbf{I},
\end{equation*}
то есть $\mathbf{Q}$ ортогональна. Подставляя $\bar{\mathbf{R}}'$ из пункта~(а) и вставляя $\bar{\mathbf{R}}^{1/2}\bar{\mathbf{R}}^{-1/2} = \mathbf{I}$ по обе стороны от $\mathbf{R}_i$, получаем
\begin{equation*}
\bar{\mathbf{R}}'^{-1/2} \mathbf{R}_i' \bar{\mathbf{R}}'^{-1/2} = \mathbf{Q} \big( \bar{\mathbf{R}}^{-1/2} \mathbf{R}_i \bar{\mathbf{R}}^{-1/2} \big) \mathbf{Q}^{\top},
\end{equation*}
и по лемме~\ref{lem:cls_conj} $\mathbf{S}_i' = \mathbf{Q}\mathbf{S}_i\mathbf{Q}^{\top}$. Отображение $\mathbf{S} \mapsto \mathbf{Q}\mathbf{S}\mathbf{Q}^{\top}$ линейно на $\mathrm{Sym}(h)$ и сохраняет норму Фробениуса; в изометричных координатах леммы~\ref{lem:cls_vec} оно задается матрицей $\mathbf{O}$, сохраняющей евклидову норму, то есть ортогональной. Матрица $\mathbf{O}$ зависит от $\mathbf{A}$ и $\bar{\mathbf{R}}$, но не от~$i$.

(в) Функционал логистической регрессии с матрицей весов $\mathbf{W} \in \mathbb{R}^{L \times d}$ и вектором смещений $\mathbf{b}$
\begin{equation*}
J(\mathbf{W}, \mathbf{b}) = \frac{1}{N} \sum_{i=1}^{N} \ell\big( y_i, \mathbf{W}\mathbf{p}_i + \mathbf{b} \big) + \lambda \|\mathbf{W}\|_F^2
\end{equation*}
строго выпукл по $\mathbf{W}$ при $\lambda > 0$, поэтому минимизатор единствен. Так как $(\mathbf{W}\mathbf{O}^{\top})(\mathbf{O}\mathbf{p}_i) = \mathbf{W}\mathbf{p}_i$ и $\|\mathbf{W}\mathbf{O}^{\top}\|_F = \|\mathbf{W}\|_F$, функционал задачи на преобразованных признаках в точке $(\mathbf{W}\mathbf{O}^{\top}, \mathbf{b})$ равен исходному в точке $(\mathbf{W}, \mathbf{b})$. Биекция $\mathbf{W} \mapsto \mathbf{W}\mathbf{O}^{\top}$ переводит минимизатор в минимизатор, откуда $\hat{\mathbf{W}}' = \hat{\mathbf{W}}\mathbf{O}^{\top}$ и $\hat{\mathbf{W}}'\mathbf{p}' = \hat{\mathbf{W}}\mathbf{p}$ для любого объекта.
\end{proof}

\begin{corollary}\label{cor:cls_units}
Результат классификации не зависит от единиц измерения сигнала, от покомпонентного масштабирования вокселей (диагональные $\mathbf{A}$), от нормировочного множителя в~\eqref{eq:cls_cov} ($\mathbf{A} = c\mathbf{I}$) и от невырожденного смешивания отобранных вокселей.
\end{corollary}

\begin{remark}
Предложение предполагает, что признаки подаются в классификатор без покоординатной стандартизации: она не коммутирует с поворотом~$\mathbf{O}$. При $h < \mathcal{T} - 1$ выборочная ковариация~\eqref{eq:cls_cov} невырождена почти наверное, и регуляризация ковариаций сжатием не требуется.
\end{remark}

\textbf{Искажение расстояний касательной проекцией.} Классификатор работает не с геодезическими расстояниями на многообразии, а с евклидовыми расстояниями между векторами~$\mathbf{p}_i$. Следующая лемма содержит ключевое вычисление: как дифференциал экспоненты меняет длину касательного вектора в метрике~\eqref{eq:cls_metric}.

\begin{lemma}\label{lem:cls_dexp}
Пусть $\mathbf{S}, \mathbf{X} \in \mathrm{Sym}(h)$, $\lambda_1, \ldots, \lambda_h$~--- собственные значения $\mathbf{S}$, а $X_{kl}$~--- элементы $\mathbf{X}$ в ортонормированном собственном базисе $\mathbf{S}$. Тогда
\begin{equation}\label{eq:cls_dexp}
\big\| d\exp_{\mathbf{S}}[\mathbf{X}] \big\|_{e^{\mathbf{S}}}^2 = \sum_{k,l=1}^{h} g(\lambda_k - \lambda_l)^2 X_{kl}^2,
\qquad
g(x) = \frac{\sinh(x/2)}{x/2}, \quad g(0) = 1 .
\end{equation}
Функция $g$ четна, $g(x) \geqslant 1$ и возрастает по $|x|$.
\end{lemma}
\begin{proof}
Касательный вектор~\eqref{eq:cls_duhamel} приложен в точке $\mathbf{P} = e^{\mathbf{S}}$, поэтому по~\eqref{eq:cls_metric} его длина равна норме Фробениуса матрицы
\begin{equation*}
e^{-\mathbf{S}/2} \Big( \int_0^1 e^{(1-u)\mathbf{S}} \mathbf{X} e^{u\mathbf{S}} du \Big) e^{-\mathbf{S}/2} = \int_0^1 e^{(1/2-u)\mathbf{S}} \mathbf{X} e^{(u-1/2)\mathbf{S}} du .
\end{equation*}
Перейдем в ортонормированный собственный базис $\mathbf{S}$; норма Фробениуса при этом не меняется, а матрицы $e^{a\mathbf{S}}$ диагональны, поэтому элемент $(k,l)$ подынтегральной матрицы равен $X_{kl} e^{(1/2-u)(\lambda_k - \lambda_l)}$. Интегрируя,
\begin{equation*}
\int_0^1 e^{(1/2-u)(\lambda_k - \lambda_l)} du = \int_{-1/2}^{1/2} e^{v(\lambda_k - \lambda_l)} dv = \frac{e^{(\lambda_k-\lambda_l)/2} - e^{-(\lambda_k-\lambda_l)/2}}{\lambda_k - \lambda_l} = g(\lambda_k - \lambda_l),
\end{equation*}
что дает~\eqref{eq:cls_dexp}. Четность $g$ очевидна, $g \geqslant 1$ следует из $\sinh x \geqslant x$ при $x \geqslant 0$, а возрастание по $|x|$~--- из возрастания $\sinh x / x$ на $[0, \infty)$.
\end{proof}

\begin{theorem}\label{thm:cls_isometry}
Для признаков~\eqref{eq:cls_feat} справедливы следующие утверждения.
\begin{enumerate}
    \item[(а)] Норма признака равна геодезическому расстоянию до опорной точки: $\|\mathbf{p}_i\|_2 = \delta_{\mathcal{R}}(\bar{\mathbf{R}}, \mathbf{R}_i)$.
    \item[(б)] Проекция не увеличивает расстояний: $\|\mathbf{p}_i - \mathbf{p}_j\|_2 \leqslant \delta_{\mathcal{R}}(\mathbf{R}_i, \mathbf{R}_j)$ для любых $i, j$.
    \item[(в)] Если $\delta_{\mathcal{R}}(\bar{\mathbf{R}}, \mathbf{R}_i) \leqslant \rho$ для всех $i$, то
    \begin{equation}\label{eq:cls_isometry}
    \delta_{\mathcal{R}}(\mathbf{R}_i, \mathbf{R}_j) \leqslant c(\rho)\, \|\mathbf{p}_i - \mathbf{p}_j\|_2,
    \qquad
    c(\rho) = \frac{\sinh(\rho/\sqrt{2})}{\rho/\sqrt{2}} .
    \end{equation}
\end{enumerate}
\end{theorem}
\begin{proof}
По предложению~\ref{prop:cls_affine} с $\mathbf{A} = \bar{\mathbf{R}}^{-1/2}$ и лемме~\ref{lem:cls_inv} обе части всех трех соотношений не меняются при переходе к матрицам $\bar{\mathbf{R}}^{-1/2}\mathbf{R}_i\bar{\mathbf{R}}^{-1/2}$, поэтому достаточно рассмотреть случай $\bar{\mathbf{R}} = \mathbf{I}$. Тогда $\mathbf{S}_i = \log \mathbf{R}_i$ и $\mathbf{R}_i = \exp \mathbf{S}_i$.

(а) По~\eqref{eq:cls_dist_log} и лемме~\ref{lem:cls_vec} $\delta_{\mathcal{R}}(\mathbf{I}, \mathbf{R}_i) = \|\log \mathbf{R}_i\|_F = \|\mathbf{S}_i\|_F = \|\mathbf{p}_i\|_2$.

(б) Из леммы~\ref{lem:cls_dexp} и $g \geqslant 1$ следует $\|d\exp_{\mathbf{S}}[\mathbf{X}]\|_{e^{\mathbf{S}}} \geqslant \|\mathbf{X}\|_F$ для всех $\mathbf{S}, \mathbf{X}$: экспонента не укорачивает касательных векторов. Пусть $\sigma$~--- произвольная гладкая кривая в $\mathbb{S}_{+}^{h}$ из $\mathbf{R}_i$ в $\mathbf{R}_j$. Поскольку $\exp$~--- диффеоморфизм, кривая $\gamma = \log \circ\, \sigma$ лежит в $\mathrm{Sym}(h)$, соединяет $\mathbf{S}_i$ с $\mathbf{S}_j$, и $\dot{\sigma}(t) = d\exp_{\gamma(t)}[\dot{\gamma}(t)]$. Для длин кривых
\begin{equation*}
L(\sigma) = \int \big\| d\exp_{\gamma(t)}[\dot{\gamma}(t)] \big\|_{\sigma(t)} dt \geqslant \int \|\dot{\gamma}(t)\|_F \, dt \geqslant \|\mathbf{S}_i - \mathbf{S}_j\|_F ,
\end{equation*}
где последнее неравенство~--- евклидов отрезок короче любой соединяющей кривой. Геодезическое расстояние есть нижняя грань $L(\sigma)$ по всем кривым, поэтому $\delta_{\mathcal{R}}(\mathbf{R}_i, \mathbf{R}_j) \geqslant \|\mathbf{S}_i - \mathbf{S}_j\|_F = \|\mathbf{p}_i - \mathbf{p}_j\|_2$.

(в) Рассмотрим отрезок $\gamma(t) = (1-t)\mathbf{S}_i + t\mathbf{S}_j$, $t \in [0,1]$, и его образ $\exp \circ\, \gamma$, соединяющий $\mathbf{R}_i$ с $\mathbf{R}_j$. По выпуклости нормы $\|\gamma(t)\|_F \leqslant \max\{\|\mathbf{S}_i\|_F, \|\mathbf{S}_j\|_F\} \leqslant \rho$, где второе неравенство~--- пункт~(а) и условие теоремы. Для симметричной $\mathbf{S}$ с $\|\mathbf{S}\|_F \leqslant \rho$ разность любых двух собственных значений оценивается как
\begin{equation*}
|\lambda_k - \lambda_l| \leqslant \sqrt{2(\lambda_k^2 + \lambda_l^2)} \leqslant \sqrt{2}\, \|\mathbf{S}\|_F \leqslant \sqrt{2}\, \rho ,
\end{equation*}
поскольку $(a-b)^2 \leqslant 2(a^2 + b^2)$ и $\sum_k \lambda_k^2 = \|\mathbf{S}\|_F^2$. Так как $g$ возрастает по модулю аргумента, из~\eqref{eq:cls_dexp} для всех $t$
\begin{equation*}
\big\| d\exp_{\gamma(t)}[\mathbf{X}] \big\|_{e^{\gamma(t)}} \leqslant g\big(\sqrt{2}\rho\big) \|\mathbf{X}\|_F = c(\rho) \|\mathbf{X}\|_F .
\end{equation*}
Следовательно, длина кривой $\exp \circ\, \gamma$ не превосходит $c(\rho) \int_0^1 \|\dot{\gamma}(t)\|_F dt = c(\rho) \|\mathbf{S}_j - \mathbf{S}_i\|_F$, а геодезическое расстояние не превосходит длины любой соединяющей кривой.
\end{proof}

\begin{remark}
Теорема~\ref{thm:cls_isometry} означает, что для всех пар объектов отношение $\delta_{\mathcal{R}}(\mathbf{R}_i, \mathbf{R}_j) / \|\mathbf{p}_i - \mathbf{p}_j\|_2$ лежит в отрезке $[1, c(\rho)]$, причем $c(\rho) \to 1$ при $\rho \to 0$. Концы отрезка отвечают двум предельным направлениям смещения в касательном пространстве (рис.~\ref{fig:cls_distortion}). Радиальные расстояния, измеряемые вдоль геодезической через опорную точку, проекция сохраняет точно: если $\mathbf{p}_k = -\mathbf{p}_i$, то касательные векторы $\mathbf{S}_i$ и $\mathbf{S}_k$ коммутируют, кривая $t \mapsto \operatorname{Exp}_{\bar{\mathbf{R}}}(t\mathbf{P}_i)$, $t \in [-1, 1]$, является геодезической, и $\delta_{\mathcal{R}}(\mathbf{R}_i, \mathbf{R}_k) = 2\rho = \|\mathbf{p}_i - \mathbf{p}_k\|_2$. Более общо, равенство в пункте~(б) выполнено для любых коммутирующих $\mathbf{S}_i$, $\mathbf{S}_j$, поскольку в общем собственном базисе $\mathbf{R}_i^{-1/2}\mathbf{R}_j\mathbf{R}_i^{-1/2} = \exp(\mathbf{S}_j - \mathbf{S}_i)$. Верхняя граница достигается на тангенциальных смещениях: для пары касательных векторов с собственными значениями $\pm \rho/\sqrt{2}$, отличающихся малым поворотом собственного базиса, отношение стремится к $c(\rho)$, так как разность собственных значений $\sqrt{2}\rho$ максимальна при заданной норме, а смещение направлено вдоль внедиагональных элементов, где множитель~\eqref{eq:cls_dexp} максимален. Таким образом, искажение порождается не удаленностью объектов друг от друга, а некоммутативностью их касательных векторов: чем сильнее различаются собственные базисы ковариационных матриц при близких спектрах, тем больше геодезическое расстояние превышает евклидово. Для облаков с невырожденным спектром касательных векторов искажение существенно меньше границы. Пункт~(б) есть частный случай общего свойства многообразий неположительной кривизны: экспоненциальное отображение не уменьшает расстояний~\cite{bhatia2007positive}. На сфере, многообразии положительной кривизны, картина обратная: радиальные расстояния через полюс также сохраняются, а тангенциальные при экспоненциальном отображении сокращаются с множителем $\sin\rho/\rho$.
\end{remark}

\begin{figure}[h!]
    \centering
    \includegraphics[width=0.62\textwidth]{cls/distortion_scheme.pdf}
    \caption{Искажение расстояний при проекции на касательное пространство. Точки $\mathbf{R}_i$, $\mathbf{R}_j$, $\mathbf{R}_k$ лежат на окружности геодезического радиуса $\rho$ вокруг $\bar{\mathbf{R}}$, их образы $\mathbf{p}_i$, $\mathbf{p}_j$, $\mathbf{p}_k$~--- на окружности того же радиуса в $T_{\bar{\mathbf{R}}}$. Радиальное расстояние через $\bar{\mathbf{R}}$ сохраняется: $\delta_{\mathcal{R}}(\mathbf{R}_i, \mathbf{R}_k) = \|\mathbf{p}_i - \mathbf{p}_k\|_2 = 2\rho$. Тангенциальное расстояние между близкими точками окружности при возвращении на многообразие растягивается; в предельном случае теоремы~\ref{thm:cls_isometry} $\delta_{\mathcal{R}}(\mathbf{R}_i, \mathbf{R}_j) = c(\rho)\,\|\mathbf{p}_i - \mathbf{p}_j\|_2$}
    \label{fig:cls_distortion}
\end{figure}

\textbf{Шум оценивания ковариаций.} Ковариационная матрица объекта оценивается по $\mathcal{T}$ сканам, поэтому часть разброса признаков есть не различие объектов, а ошибка оценивания. Следующая теорема показывает, что распределение этой ошибки в геодезической метрике не зависит от истинной ковариации, и дает ее оценку. Используется классический результат о сингулярных числах гауссовой матрицы.

\begin{proposition}[{\cite[следствие~5.35]{vershynin2012introduction}}]\label{prop:cls_vershynin}
Пусть $\mathbf{G} \in \mathbb{R}^{n \times k}$, $n \geqslant k$, имеет независимые элементы с распределением $\mathcal{N}(0,1)$. Тогда для любого $u \geqslant 0$ с вероятностью не меньше $1 - 2e^{-u^2/2}$
\begin{equation}\label{eq:cls_vershynin}
\sqrt{n} - \sqrt{k} - u \leqslant s_{\min}(\mathbf{G}) \leqslant s_{\max}(\mathbf{G}) \leqslant \sqrt{n} + \sqrt{k} + u ,
\end{equation}
где $s_{\min}$ и $s_{\max}$~--- наименьшее и наибольшее сингулярные числа.
\end{proposition}

\begin{theorem}\label{thm:cls_conc}
Пусть столбцы матрицы $\mathbf{Y} \in \mathbb{R}^{h \times \mathcal{T}}$~--- независимые случайные векторы с распределением $\mathcal{N}(\mathbf{0}, \mathbf{R})$, $\mathbf{R} \in \mathbb{S}_{+}^{h}$, и $\hat{\mathbf{R}} = \mathcal{T}^{-1} \mathbf{Y}\mathbf{Y}^{\top}$. Тогда
\begin{enumerate}
    \item[(а)] распределение случайной величины $\delta_{\mathcal{R}}(\hat{\mathbf{R}}, \mathbf{R})$ не зависит от $\mathbf{R}$, а только от $h$ и $\mathcal{T}$;
    \item[(б)] для любого $t > 0$ такого, что $\varepsilon = \sqrt{h/\mathcal{T}} + t < 1$, с вероятностью не меньше $1 - 2\exp(-\mathcal{T}t^2/2)$
    \begin{equation}\label{eq:cls_conc}
    \delta_{\mathcal{R}}(\hat{\mathbf{R}}, \mathbf{R}) \leqslant \frac{2\sqrt{h}\, \varepsilon}{1 - \varepsilon} .
    \end{equation}
\end{enumerate}
\end{theorem}
\begin{proof}
(а) Положим $\boldsymbol{\Xi} = \mathbf{R}^{-1/2}\mathbf{Y}$. Столбцы $\boldsymbol{\Xi}$ независимы и распределены как $\mathcal{N}(\mathbf{0}, \mathbf{I})$, то есть все элементы $\boldsymbol{\Xi}$ независимы и стандартны нормальны, и их распределение от $\mathbf{R}$ не зависит. По лемме~\ref{lem:cls_inv} с $\mathbf{A} = \mathbf{R}^{-1/2}$
\begin{equation*}
\delta_{\mathcal{R}}(\hat{\mathbf{R}}, \mathbf{R}) = \delta_{\mathcal{R}}\big( \mathbf{R}^{-1/2}\hat{\mathbf{R}}\mathbf{R}^{-1/2}, \mathbf{I} \big) = \big\| \log \hat{\mathbf{W}} \big\|_F ,
\qquad
\hat{\mathbf{W}} = \mathcal{T}^{-1} \boldsymbol{\Xi}\boldsymbol{\Xi}^{\top},
\end{equation*}
и правая часть есть функция только от $\boldsymbol{\Xi}$.

(б) Применим предложение~\ref{prop:cls_vershynin} к матрице $\boldsymbol{\Xi}^{\top} \in \mathbb{R}^{\mathcal{T} \times h}$ с $u = t\sqrt{\mathcal{T}}$: с вероятностью не меньше $1 - 2\exp(-\mathcal{T}t^2/2)$ все ее сингулярные числа лежат в отрезке $\sqrt{\mathcal{T}}\,[1 - \varepsilon, 1 + \varepsilon]$. Собственные значения $\hat{\mathbf{W}}$ суть квадраты сингулярных чисел $\boldsymbol{\Xi}^{\top}$, деленные на $\mathcal{T}$, поэтому на этом событии они лежат в $[(1-\varepsilon)^2, (1+\varepsilon)^2]$, и для каждого из них
\begin{equation*}
\big| \log \lambda_k(\hat{\mathbf{W}}) \big| \leqslant \max\big\{ 2\log(1+\varepsilon),\ -2\log(1-\varepsilon) \big\} = -2\log(1-\varepsilon) \leqslant \frac{2\varepsilon}{1-\varepsilon} ,
\end{equation*}
где равенство следует из $-\log(1-\varepsilon) \geqslant \log(1+\varepsilon)$ при $\varepsilon \in (0,1)$, а последнее неравенство~--- из $\log x \leqslant x - 1$ при $x = 1/(1-\varepsilon)$. Остается заметить, что $\|\log\hat{\mathbf{W}}\|_F^2 = \sum_{k=1}^{h} \log^2 \lambda_k(\hat{\mathbf{W}}) \leqslant h \max_k \log^2 \lambda_k(\hat{\mathbf{W}})$, что дает~\eqref{eq:cls_conc}.
\end{proof}

\begin{corollary}\label{cor:cls_stab}
Пусть признаки $\mathbf{p}_i$ и $\hat{\mathbf{p}}_i$ вычислены по~\eqref{eq:cls_feat} при одной и той же опорной точке по истинной ковариации $\mathbf{R}_i$ и по ее оценке $\hat{\mathbf{R}}_i$. Тогда $\|\hat{\mathbf{p}}_i - \mathbf{p}_i\|_2 \leqslant \delta_{\mathcal{R}}(\hat{\mathbf{R}}_i, \mathbf{R}_i)$, и эта величина ограничена теоремой~\ref{thm:cls_conc}.
\end{corollary}
\begin{proof}
Первое неравенство~--- пункт~(б) теоремы~\ref{thm:cls_isometry}, примененный к паре $\hat{\mathbf{R}}_i, \mathbf{R}_i$.
\end{proof}

\begin{remark}
Пункт~(а) теоремы~\ref{thm:cls_conc} означает, что типичную величину ошибки оценивания ковариации в геодезической метрике можно получить моделированием с единичной ковариацией, и результат будет верен для любых реальных ковариаций с теми же $h$ и $\mathcal{T}$. Сравнение этой величины с наблюдаемым разбросом $\delta_{\mathcal{R}}(\bar{\mathbf{R}}, \mathbf{R}_i)$ показывает, какая доля разброса признаков обусловлена шумом оценивания, а не различием объектов. Оценка~\eqref{eq:cls_conc} задает скорость $\delta_{\mathcal{R}}(\hat{\mathbf{R}}, \mathbf{R}) = O\big(h/\sqrt{\mathcal{T}}\big)$: шум оценивания растет линейно по числу отобранных вокселей и убывает как корень из длины ряда. Отсюда следует ограничение $h \ll \mathcal{T}$ на параметр $h$ Алгоритма~\ref{alg:bad}, согласованное с выбором в разделе~\ref{sec4:cls}. Для центрированных рядов и нормировки~\eqref{eq:cls_cov} теорема применяется с $\mathcal{T} - 1$ вместо $\mathcal{T}$; выбор нормировочного множителя на признаки не влияет по следствию~\ref{cor:cls_units}.
\end{remark}

\section{Вычислительный эксперимент}\label{sec4:cls}

Вычислительный эксперимент использует данные фМРТ из исследования~\cite{haxby2001distributed}, собранные при изучении представлений лиц и объектов в зрительной коре. Набор данных содержит записи фМРТ от 6 испытуемых. Каждый испытуемый проходил 12 сессий. В каждой сессии испытуемый пассивно просматривал черно-белые изображения восьми категорий, сгруппированные в блоки по 24 секунды, разделенные периодами отдыха. Каждое изображение демонстрировалось в течение 500 мс, после чего следовал интервал между стимулами длительностью 1500 мс. Данные фМРТ мозга записывались с частотой 2,5~Гц. Примеры стимулов приведены на рис.~\ref{fig:cls_stimuli}, характеристики записей~--- в Таблице~\ref{table:cls_parameters}.
\begin{figure}[h!]
    \centering
    \includegraphics[width=\textwidth]{figs/cls_stimuli.pdf}
    \caption{Визуализация стимулов из используемого набора данных. В наборе данных представлено 8 различных типов визуальных стимулов. Для каждого типа выбрано по 4 примера изображений. Стимулы для каждой категории значительно различаются.}
    \label{fig:cls_stimuli}
\end{figure}

\begin{table}[h!]
    \centering
    \caption{Основные параметры временных рядов фМРТ из выборки}\label{table:cls_parameters}
    \begin{tabular}{lcc}
        \toprule
        Название & Обозначение & Значение \\
        \midrule
        Количество сканов за сессию & $\mathcal{T}_{\mathrm{rec}}$ & 121 \\
        Частота сканирования & $\mu$ & 2,5 Гц \\
        Размеры скана & $X, Y, Z$ & 40, 64, 64 \\
        \bottomrule
    \end{tabular}
\end{table}

Кроме того, набор данных содержит временные ряды стимулов и маски активных областей мозга, полученные нейроучеными. Эти области были определены для каждого участника эксперимента индивидуально. В наборе данных Haxby et al.~\cite{haxby2001distributed} учитывается гемодинамический лаг в процессе аннотирования, что обеспечивает корректную временную синхронизацию данных. Это позволяет установить гиперпараметр $\Delta t$ равным $0$ в ходе экспериментов.

\subsection{Аугментация данных}\label{sec4:cls_aug}

Для формирования обучающей выборки к записи каждого испытуемого применялся метод сегментации из раздела~\ref{sec3:cls_aug}. В Таблице~\ref{table:cls_augmentation} представлены параметры выборок, полученных в результате аугментации: из записи каждого испытуемого получено $N = 96$ временных рядов длины $\mathcal{T} = 19$ по $L = 8$ категориям, по $N_y = 12$ объектов на категорию. Выборка каждого испытуемого разбита на обучающую и тестовую в пропорции 80\% к 20\%. Качество оценивается метриками Accuracy, Macro F1-score и Micro F1-score.

\begin{table}[h!]
    \centering
    \caption{Параметры выборки каждого испытуемого, полученной в результате аугментации временных рядов фМРТ}\label{table:cls_augmentation}
    \begin{tabular}{lcc}
        \toprule
        Название & Обозначение & Значение \\
        \midrule
        Количество объектов & $N$ & 96 \\
        Длина каждого сегмента & $\mathcal{T}$ & 19 \\
        Количество классов стимулов & $L$ & 8 \\
        Количество объектов на класс & $N_y$ & 12 \\
        \bottomrule
    \end{tabular}
\end{table}

\subsection{Оценка алгоритма выделения масок}\label{sec4:cls_bad}

Для оценки Алгоритма~\ref{alg:bad} его маски активности сравнивались с целевыми масками нейроученых и масками, полученными в ходе статистического анализа, с использованием метрик precision, recall и Intersection over Union (IoU), усредненных по шести испытуемым. Precision измеряет долю правильных положительных предсказаний алгоритма, recall количественно оценивает способность алгоритма идентифицировать все релевантные области, а IoU оценивает перекрытие между предсказанными и истинными масками.

Гиперпараметры $k_s$ и $h$ выбирались для максимизации recall, чтобы маски активности покрывали как можно больше аннотаций нейроученых, даже если это означало включение некоторых дополнительных областей. Такой подход максимизирует выявление всех релевантных активных областей, потенциально за счет включения некоторых нерелевантных. Кроме того, важно было поддерживать высокую точность относительно статистически значимой маски. По результатам предварительных экспериментов значения параметров были установлены $k_s = 4$ и $h = 10$.

\begin{table}[h!]
    \centering
    \caption{Усредненные значения метрик для сравнения масок активности. В таблице представлены средние значения и стандартные отклонения для каждой метрики}\label{table:cls_mask_metrics}
    \begin{tabular}{lccc}
        \toprule
        Метрика & Precision & Recall & IoU \\
        \midrule
        Алгоритм~\ref{alg:bad} vs. Целевая & 0.11 $\pm$ 0.04 & 0.43 $\pm$ 0.07 & 0.10 $\pm$ 0.04 \\
        Алгоритм~\ref{alg:bad} vs. Статистически значимая & 1.00 $\pm$ 0.00 & 0.22 $\pm$ 0.09 & 0.22 $\pm$ 0.09 \\
        \bottomrule
    \end{tabular}
\end{table}

Результаты, представленные в Таблице~\ref{table:cls_mask_metrics}, показывают, что алгоритм эффективно выявляет области мозга со статистически значимой корреляцией со стимулами, хорошо согласуясь с аннотациями нейроученых.

Для визуальной демонстрации рассматривается выборка первого испытуемого с размером ядра $k_s = 4$ и $h = 10$. Выбраны срезы на различных координатах: 20-й, 28-й и 24-й срезы вдоль каждой из трех осей. На рис.~\ref{fig:cls_masks} приведены срезы скана из выборки с соответствующими масками активности: в верхней строке~--- области, определенные Алгоритмом~\ref{alg:bad}, в средней~--- области, полученные в результате статистического анализа, в нижней~--- аннотации нейроученых, служащие эталоном. Экспериментально подтверждено, что корреляция между взвешенными областями и стимулами статистически значима, а выявленные области хорошо согласуются с аннотациями нейроученых.

\begin{figure}[h!]
    \centering
    \begin{tabular}{ccc}
        \begin{subfigure}{0.3\textwidth}
            \centering
            \includegraphics[height=1.\textwidth]{cls/Figure_5.png}
            \caption{Алгоритм~\ref{alg:bad}}
        \end{subfigure}
        &
        \begin{subfigure}{0.3\textwidth}
            \centering
            \includegraphics[height=1.\textwidth]{cls/Figure_6.png}
            \caption{Алгоритм~\ref{alg:bad}}
        \end{subfigure}
        &
        \begin{subfigure}{0.3\textwidth}
            \centering
            \includegraphics[height=1.\textwidth]{cls/Figure_7.png}
            \caption{Алгоритм~\ref{alg:bad}}
        \end{subfigure}
        \\
        \begin{subfigure}{0.3\textwidth}
            \centering
            \includegraphics[height=1.\textwidth]{cls/Figure_8.png}
            \caption{Статистически значимые}
        \end{subfigure}
        &
        \begin{subfigure}{0.3\textwidth}
            \centering
            \includegraphics[height=1.\textwidth]{cls/Figure_9.png}
            \caption{Статистически значимые}
        \end{subfigure}
        &
        \begin{subfigure}{0.3\textwidth}
            \centering
            
\includegraphics[height=1.\textwidth]{cls/Figure_10.png}
            \caption{Статистически значимые}
        \end{subfigure}
        \\
        \begin{subfigure}{0.3\textwidth}
            \centering
            \includegraphics[height=1.\textwidth]{cls/Figure_11.png}
            \caption{Аннотации нейроученых}
        \end{subfigure}
        &
        \begin{subfigure}{0.3\textwidth}
            \centering
            \includegraphics[height=1.\textwidth]{cls/Figure_12.png}
            \caption{Аннотации нейроученых}
        \end{subfigure}
        &
        \begin{subfigure}{0.3\textwidth}
            \centering
            \includegraphics[height=1.\textwidth]{cls/Figure_13.png}
            \caption{Аннотации нейроученых}
        \end{subfigure}
    \end{tabular}
    \caption{Визуализация масок активности фМРТ. Выбранные области обозначены красным цветом. На верхней строке показаны предсказанные маски активности с использованием Алгоритма~\ref{alg:bad}, на средней строке~--- статистически значимые области, полученные в результате анализа, а на нижней строке~--- аннотации нейроученых в качестве эталона. Предсказанные области достаточно покрывают истинные аннотации, а статистически значимая корреляция взвешенных вокселей со стимулами подчеркивает надежность алгоритма}
    \label{fig:cls_masks}
\end{figure}

Также рассматривалось качество работы алгоритма на неинформативных данных. Использовались временные ряды фМРТ первого испытуемого и случайно сгенерированный бинарный временной ряд стимулов. На рис.~\ref{fig:cls_correctness} показаны срезы скана фМРТ, где цветом выделены области мозга, предсказанные Алгоритмом~\ref{alg:bad}: алгоритм предсказывает случайные области, не связанные с активностью мозга. Как и в предыдущем эксперименте, были проведены статистические тесты. Ни одна из полученных областей не оказалась статистически значимой, что подтверждает корректность алгоритма.

Высокая пикселизация на рис.~\ref{fig:cls_masks} и~\ref{fig:cls_correctness} является результатом намеренного сжатия сканов фМРТ в Алгоритме~\ref{alg:bad} для устранения лишнего шума. Визуализированные маски активности показаны после увеличения разрешения ($\mathcal{M}^{\uparrow}$), однако модель классификации работает со сжатыми сканами и маской $\mathcal{M}$, минуя шаг увеличения разрешения.

\begin{figure}[h!]
    \centering
    \includegraphics[height=0.3\textwidth]{cls/Figure_14.png}
    \hfill
    \includegraphics[height=0.3\textwidth]{cls/Figure_15.png}
    \hfill
    \includegraphics[height=0.3\textwidth]{cls/Figure_16.png}
    \caption{\textbf{Анализ работы Алгоритма~\ref{alg:bad} на неинформативных данных.} Представлены срезы по трем измерениям. Использовался случайно сгенерированный бинарный временной ряд стимулов. Алгоритм предсказывает случайные области, не связанные с активностью мозга}
    \label{fig:cls_correctness}
\end{figure}

\subsection{Результаты классификации}\label{sec4:cls_results}

Рассматривались данные фМРТ первых трех испытуемых. В ходе экспериментов были зафиксированы следующие параметры модели: количество областей для каждого класса $h_y = 10$, $y = 1, \ldots, L$, и размер ядра $k_s = 4$. В качестве классификатора $\boldsymbol{\eta}$ использовалась логистическая регрессия со стратегией <<один против всех>> и единичным коэффициентом $L_2$-регуляризации. Кроме того, рассматривался двухслойный перцептрон, каждый слой которого содержал 128 нейронов с сигмоидной функцией активации.

В Таблице~\ref{table:cls_results} представлены усредненные значения метрик по испытуемым на тестовой выборке, а также их стандартные отклонения и 95\% доверительные интервалы (англ. Confidence Interval, CI), которые рассчитаны для определения диапазона, в котором, вероятно, находятся истинные значения метрик.

\begin{table}[h!]
    \centering
    \caption{Средние значения метрик по испытуемым для тестовой выборки, включая стандартные отклонения и 95\% доверительные интервалы (англ. Confidence Interval, CI)}\label{table:cls_results}
    \resizebox{\textwidth}{!}{
        \begin{tabular}{lcccccc}
            \toprule
            Классификатор $\boldsymbol{\eta}$ & Accuracy & Macro F1-Score & Micro F1-Score & Accuracy CI & Macro F1 CI & Micro F1 CI \\
            \midrule
            Логистическая регрессия & 0.61 $\pm$ 0.04 & 0.56 $\pm$ 0.03 & 0.61 $\pm$ 0.04 & [0.56, 0.66] & [0.52, 0.60] & [0.56, 0.66] \\
            Перцептрон & 0.64 $\pm$ 0.05 & 0.60 $\pm$ 0.05 & 0.62 $\pm$ 0.04 & [0.57, 0.71] & [0.54, 0.67] & [0.57, 0.67] \\
            \bottomrule
        \end{tabular}
    }
\end{table}

Полученные метрики указывают на возможность улучшения, что ожидаемо при ограниченном объеме данных. Тем не менее, модель демонстрирует перспективы в обучении зависимостям даже при малом объеме данных и значительном количестве классов. Для дальнейшего анализа исследовалось влияние доли обучающей выборки на качество; метрики последовательно рассчитывались на одной и той же тестовой выборке для обеспечения сопоставимости.

На рис.~\ref{fig:cls_fraction} представлены графики Accuracy и Macro F1-score для логистической регрессии и перцептрона в зависимости от доли обучающей выборки, изменяющейся от 0 до 1, с погрешностями в виде стандартного отклонения по испытуемым. Устойчивый восходящий тренд указывает на то, что извлеченные пространственно-временные характеристики значимы и что качество модели улучшается с увеличением объема обучающих данных: улучшения не являются случайными или следствием переобучения, а отражают реальное повышение способности модели обучаться на данных.

\begin{figure}[h!]
    \centering
    \begin{subfigure}{0.49\textwidth}
        \centering
        \includegraphics[width=\textwidth]{cls/Figure_17.pdf}
        \caption{Accuracy}
    \end{subfigure}
    \hfill
    \begin{subfigure}{0.49\textwidth}
        \centering
        \includegraphics[width=\textwidth]{cls/Figure_18.pdf}
        \caption{Macro F1-Score}
    \end{subfigure}
    \caption{Анализ улучшения качества декодирования в зависимости от доли обучающей выборки. Графики представляют метрики Accuracy и Macro F1-score, усредненные по испытуемым, в зависимости от доли обучающей выборки, изменяющейся от 0 до 1. Метрики вычислялись на одной и той же тестовой выборке для обеспечения согласованности. Результаты демонстрируют последовательное улучшение качества декодирования по мере увеличения размера обучающей выборки}
    \label{fig:cls_fraction}
\end{figure}

\subsection{Анализ исключением компонент}\label{sec4:cls_ablation}
Для оценки вклада каждого компонента в предложенную модель было проведено исследование методом исключения компонент. Рассматривались следующие варианты моделей:
\begin{itemize}
    \item \textit{Без TSM}: временные ряды преобразуются в вектор после применения масок каждого класса, конкатенируются и подаются на вход классификатору без использования TSM.
    \item \textit{Без масок классов}: маска активности рассчитывается как среднее по всем временным рядам, без учета специфичных для классов масок, и TSM применяется для кодирования.
    \item \textit{Оригинальные маски}: вместо масок, сгенерированных алгоритмом из раздела~\ref{sec3:cls_mask}, используются оригинальные маски, предоставленные нейроучеными, и TSM применяется для кодирования временных рядов.
\end{itemize}

В Таблице~\ref{table:ablation} представлены усредненные значения метрик по испытуемым, а также их стандартные отклонения и 95\% доверительные интервалы. Статистическая значимость оценивалась с использованием t-тестов для связанных выборок, с результатами, обозначенными звездочками, где *$p\text{-value}<.05$ и **$p\text{-value}<.005$. Это подчеркивает важность каждого компонента относительно полной модели. Была применена поправка Холма для множественного тестирования, чтобы обеспечить достоверность статистических сравнений по трем метрикам. В качестве классификатора $ \boldsymbol{\eta}$ использовалась логистическая регрессия.

\begin{table}[h!]
    \centering
    \caption{Анализ методом исключения компонент влияния отдельных компонентов модели на качество классификации. Метрики на тестовых выборках усреднены по испытуемым. Статистическая значимость обозначена звездочками: *$p\text{-value} < .05$, **$p\text{-value} < .005$, отражающими значимость улучшений относительно полной модели. Доверительные интервалы (англ. Confidence Interval, CI) представляют диапазон, в котором, вероятно, находятся истинные значения метрик с 95\% доверительной вероятностью.}\label{table:ablation}
    \resizebox{\textwidth}{!}{
        \begin{tabular}{lccccccccc}
            \toprule
            Метод & Accuracy & Macro F1-Score & Micro F1-Score & Accuracy CI & Macro F1 CI & Micro F1 CI \\
            \midrule
            Без TSM & 0.24 $\pm$ 0.11* & 0.22 $\pm$ 0.11* & 0.24 $\pm$ 0.11* & [0.11, 0.37] & [0.09, 0.36] & [0.10, 0.37] \\
            Без масок классов & 0.25 $\pm$ 0.13* & 0.23 $\pm$ 0.13* & 0.25 $\pm$ 0.13* & [0.09, 0.41] & [0.07, 0.39] & [0.09, 0.41] \\
            Оригинальные маски & 0.10 $\pm$ 0.05** & 0.09 $\pm$ 0.04** & 0.10 $\pm$ 0.05** & [0.04, 0.16] & [0.03, 0.14] & [0.04, 0.16] \\
            Предложенная модель & \textbf{0.61 $\pm$ 0.04} & \textbf{0.56 $\pm$ 0.03} & \textbf{0.61 $\pm$ 0.04} & [0.56, 0.66] & [0.52, 0.60] & [0.56, 0.66] \\
            \bottomrule
        \end{tabular}
    }
\end{table}

Исследование методом исключения компонент показывает, что метрики значительно снижаются при исключении проецирования на касательное пространство риманова многообразия, что подчеркивает критическую роль этого шага в предложенной модели. Кроме того, отсутствие масок активности для каждого класса стимулов также приводит к снижению качества классификации, демонстрируя важность учета индивидуальных характеристик каждого стимула. Статистическая значимость, обозначенная звездочками, подтверждает, что наблюдаемые снижения производительности статистически значимы.

\subsection{Сравнение с нейросетевыми подходами}\label{sec4:cls_nn}
Для демонстрации превосходства предложенной модели над подходами на основе нейронных сетей, особенно при работе с ограниченными данными, мы сравниваем ее с двумя альтернативными моделями, использующими различные методы кодирования временных рядов фМРТ:
\begin{itemize}
    \item \textit{Оригинальные маски и LSTM}: эта модель сначала применяет оригинальные маски, предоставленные нейроучеными, для выбора активных областей мозга. Затем временные ряды из этих областей кодируются с использованием сети долгой краткосрочной памяти (англ. Long Short-Term Memory, LSTM), которая улавливает временные зависимости. Извлеченные признаки передаются через линейный слой с активацией softmax для классификации.
    \item \textit{Оригинальные маски и Attention}: эта модель сначала применяет оригинальные маски, предоставленные нейроучеными, для выбора активных областей мозга. Затем временные ряды из этих областей кодируются с использованием слоя внимания для кодирования временных рядов. Механизм внимания взвешивает временные признаки, фокусируясь на наиболее релевантных временных точках, и результирующий контекстный вектор передается через линейный слой с активацией softmax для классификации.
\end{itemize}

В Таблице~\ref{table:model_comparison} представлены средние значения метрик по испытуемым на тестовой выборке для каждой модели в сравнении с предложенной моделью, использующей логистическую регрессию в качестве классификатора $\boldsymbol{\eta}$. Таблица включает статистическую значимость, оцененную с использованием t-тестов для связанных выборок, обозначенную звездочками, где **$p\text{-value} < .005$, рассчитанную для проверки гипотезы об том, что улучшения метрик для полной модели статистически значимы. Применяется поправка Холма для множественного тестирования, чтобы обеспечить достоверность статистических сравнений по трем метрикам. Кроме того, представлены 95\% доверительные интервалы (англ. Confidence Interval, CI), чтобы указать диапазон, в котором, вероятно, находятся истинные значения метрик.

\begin{table}[h!]
    \centering
    \caption{Сравнение модели с подходами на основе нейронных сетей. Метрики усреднены по испытуемым для тестовой выборки, с указанием статистической значимости и доверительных интервалов. Статистическая значимость обозначена звездочками: **$p\text{-value} < .005$.}\label{table:model_comparison}
    \resizebox{\textwidth}{!}{
        \begin{tabular}{lcccccc}
            \toprule
            Модель & Accuracy & Macro F1-Score & Micro F1-Score & Accuracy CI & Macro F1 CI & Micro F1 CI \\
            \midrule
            Оригинальные маски и LSTM & 0.12 $\pm$ 0.02** & 0.03 $\pm$ 0.01** & 0.12 $\pm$ 0.02** & [0.09, 0.14] & [0.02, 0.04] & [0.09, 0.14] \\
            Оригинальные маски и Attention & 0.14 $\pm$ 0.04** & 0.04 $\pm$ 0.03** & 0.14 $\pm$ 0.04** & [0.11, 0.18] & [0.02, 0.07] & [0.11, 0.18] \\
            Предложенная модель & \textbf{0.61 $\pm$ 0.04} & \textbf{0.56 $\pm$ 0.03} & \textbf{0.61 $\pm$ 0.04} & [0.56, 0.66] & [0.52, 0.60] & [0.56, 0.66] \\
            \bottomrule
        \end{tabular}
    }
\end{table}

Экспериментальные результаты, представленные в Таблице~\ref{table:model_comparison}, подчеркивают превосходство предложенной модели по сравнению с подходами на основе нейронных сетей, особенно в сценариях с ограниченной выборкой. Использование специфичных для классов масок из раздела~\ref{sec3:cls_mask}, и TSM для кодирования позволяет модели эффективно улавливать как пространственные, так и временные характеристики данных фМРТ. Статистическая значимость, обозначенная $p\text{-value} < .005$, подтверждает гипотезу о том, что наблюдаемые улучшения метрик статистически значимы и маловероятно произошли случайно. Доверительные интервалы дополнительно обеспечивают меру надежности этих оценок, укрепляя обоснованность улучшений производительности модели.

\section{Обсуждение результатов}\label{sec4:cls_disc}
Экспериментальные результаты показали, что исключение отдельных компонентов предложенной модели значительно снижает качество классификации. Это подчеркивает важность пространственно-временных характеристик, извлекаемых с помощью данной методологии, для достижения высокой точности классификации. Основное внимание в этой главе уделяется оценке качества классификации в условиях ограниченной выборки, где традиционные подходы демонстрируют низкую эффективность. Таким образом, предложенная методология представляет ценность для анализа и интерпретации нейронных данных, особенно в ситуациях с ограниченным объемом данных.

В сравнении с подходами на основе нейронных сетей предложенная методология показывает более высокие результаты. Нейронные сети часто требуют больших объемов данных для эффективного обучения зависимостей и паттернов. В то же время предложенный подход, использующий логистическую регрессию и перцептрон, демонстрирует лучшее качество декодирования на небольших объемах выборки. Это достигается за счет эффективного использования масок активности и отображения на касательное пространство, которые учитывают как пространственные, так и временные характеристики данных фМРТ.

Тем не менее, используемый набор данных является уникальным и содержит специфический набор временных точек стимуляции. Кроме того, небольшое количество участников может не полностью отражать разнообразие активности мозга. Валидация методологии на более крупных и разнообразных наборах данных, а также изучение альтернативных схем взвешивания для повышения производительности составляют направление дальнейшей работы. Несмотря на эти ограничения, полученные результаты способствуют лучшему пониманию реакции человеческого мозга на внешние стимулы и закладывают основу для последующих исследований.

\fixme{Дополнить главу: гарантия скрининга масок активности (рабочий черновик: drafts/theory\_classification.tex); численная проверка теорем раздела~\ref{sec3:cls_theory} на данных эксперимента (отношение геодезических и касательных расстояний, шумовой радиус).}
